\chapter{转移算子、有限条带谱与相关长度}
\label{chap:phys-sm-transfer-matrix-fermionization}

图展开把二维 Ising 的配分函数组织成边集之和；要得到完整自由能和有限尺寸谱，更有效的做法是把一个方向的统计求和改写为线性算子的幂。先建立 transfer operator 的一般有限条带结构：精确矩阵元、Perron--Frobenius 主本征值、谱表示的相关函数以及 Pauli 算符分解。随后再引入费米子化，从而把“转移算子存在”与“方格 Ising 恰好能被 Clifford 代数线性化”分开。

\section{逐行配置空间与精确转移算子}
\label{sec:phys-sm-transfer-matrix-foundations}

考虑 \(M\times N\) 各向异性方格，水平方向耦合 \(K_h=\beta J_h\)，垂直方向耦合 \(K_v=\beta J_v\)。一行配置记为
\[
 \boldsymbol\sigma=(\sigma_1,\ldots,\sigma_N),
 \qquad \sigma_j=\pm1.
\]
所有行配置张成 \(2^N\) 维配置向量空间
\begin{equation}
 \mathscr R_N=(\mathbb C^2)^{\otimes N}.
 \label{eq:phys-sm-row-space}
\end{equation}
这里的 ket 只是经典行配置的线性编码，不是额外的量子时间演化 Hilbert 空间。

横向取周期边界 \(\sigma_{N+1}=\sigma_1\)。把一行内部的水平 Boltzmann 权重平均分到相邻两次纵向传播，定义对称 transfer operator \(\mathsf T_N\) 的矩阵元
\begin{equation}
 \langle\boldsymbol\sigma'|\mathsf T_N|\boldsymbol\sigma\rangle
 =\exp\!\left[
 K_v\sum_{j=1}^{N}\sigma_j'\sigma_j
 +\frac{K_h}{2}\sum_{j=1}^{N}
 \left(\sigma_j\sigma_{j+1}
 +\sigma_j'\sigma_{j+1}'\right)
 \right].
 \label{eq:phys-sm-row-transfer-matrix}
\end{equation}
纵向也取周期边界时，有限环面配分函数严格等于
\begin{equation}
 Z_{N,M}=\Tr_{\mathscr R_N}\mathsf T_N^{\,M}.
 \label{eq:phys-sm-transfer-trace-general}
\end{equation}
这个恒等式在任意有限 \(M,N\) 成立，尚未取任何热力学极限。

\section{Perron--Frobenius 与长柱面极限}

有限实耦合下，\(\mathsf T_N\) 的每个矩阵元都严格为正。Perron--Frobenius 定理因此给出唯一的最大正本征值 \(\Lambda_0(N)>0\)，且
\[
 |\Lambda_a(N)|<\Lambda_0(N),\qquad a\ge1.
\]
固定宽度 \(N\) 后，
\begin{equation}
 \lim_{M\to\infty}\frac1M\log Z_{N,M}
 =\log\Lambda_0(N),
 \label{eq:phys-sm-transfer-long-cylinder}
\end{equation}
因此每格点有限宽自由能为
\begin{equation}
 -\beta f_N=\frac1N\log\Lambda_0(N).
 \label{eq:phys-sm-strip-free-energy}
\end{equation}
真正的二维自由能还需再取 \(N\to\infty\)。Perron--Frobenius 只控制固定 \(N\) 的有限维谱，并不能单独证明第二个极限。

\begin{derivation}[主本征值为何控制长柱面自由能]
由 Perron--Frobenius 定理，正矩阵 \(\mathsf T_N\) 有严格主本征值 \(\Lambda_0>0\)：\(\Lambda_0\) 是单根，对应严格正特征向量 \(|0\rangle\)，且所有其他本征值满足 \(|\Lambda_a|<\Lambda_0\)。写出全部本征值
\[
 \Lambda_0,\Lambda_1,\ldots,\Lambda_{2^N-1},
 \qquad |\Lambda_a|<\Lambda_0\;(a\ge1).
\]
这里需区分两种版本的 Perron--Frobenius：对\emph{严格正矩阵}（每个元素 \(>0\)）主本征值严格占优且为非退化单根；对\emph{不可约非负矩阵}（存在 \(k\) 使 \(\mathsf T^k\) 严格正）仍有正主本征值，但次主本征值可能与之等模（例如临界点上方简并）。铁磁 Ising 的 \(\mathsf T_N\) 因 Boltzmann 权 \(\ee^{\beta J\sigma\sigma'}>0\) 对所有 \(\sigma,\sigma'\) 严格为正，落入第一类，故谱隙
\[
 \Delta_N:=\log\Lambda_0-\log|\Lambda_1|>0
\]
严格为正，并直接给出沿纵向的关联长度上界 \(\xi_\parallel(N)=1/\Delta_N\)。

由\eqnrefs{eq:phys-sm-transfer-trace-general}，配分函数是本征值的 \(M\) 次幂之和：
\[
 Z_{N,M}
 =\sum_{a=0}^{2^N-1}\Lambda_a^M
 =\Lambda_0^M\left[
 1+\sum_{a=1}^{2^N-1}
 \left(\frac{\Lambda_a}{\Lambda_0}\right)^M
 \right].
\]
对每个 \(a\ge1\)，比值 \(r_a=\Lambda_a/\Lambda_0\) 满足 \(|r_a|<1\)，故 \(|r_a|^M\to0\)（\(M\to\infty\)）。更精确地，令 \(r_{\max}=\max_{a\ge1}|\Lambda_a|/\Lambda_0<1\)，则
\[
 \left|\sum_{a=1}^{2^N-1}r_a^M\right|
 \le(2^N-1)\,r_{\max}^M
 =o(1),
\]
因此方括号 \(1+o(1)\to1\)。取对数，
\[
 \frac1M\log Z_{N,M}
 =\log\Lambda_0
 +\frac1M\log\!\left[1+\sum_{a\ge1}r_a^M\right].
\]
余项满足
\[
 \left|\frac1M\log(1+o(1))\right|
 \le\frac{C\,(2^N-1)\,r_{\max}^M}{M}
 \le\frac{C\,(2^N-1)\,\ee^{-M\Delta_N}}{M}
 \longrightarrow0,
\]
其中用了 \(|\log(1+x)|\le2|x|\)（\(|x|\) 充分小），并把 \(r_{\max}^M=\ee^{-M\Delta_N}\) 写成谱隙的指数衰减形式。因此
\[
 \frac1M\log Z_{N,M}\longrightarrow\log\Lambda_0,
\]
再除以 \(N\) 即得\eqnrefs{eq:phys-sm-strip-free-energy}。主本征值的严格性来自 \(\mathsf T_N\) 的全部元素为正（Perron--Frobenius），而铁磁 Ising 的 transfer matrix 恰满足此条件。

取 \(N=1\)，单列即单自旋，\(\mathsf T_1\) 退化为 \(2\times2\) 矩阵
\begin{equation*}
 \mathsf T_1
 =\begin{pmatrix}
 \ee^{\beta J} & \ee^{-\beta J}\\
 \ee^{-\beta J} & \ee^{\beta J}
 \end{pmatrix},
\end{equation*}
本征值为 \(\Lambda_0=2\cosh(\beta J)\)、\(\Lambda_1=2\sinh(\beta J)\)，谱隙 \(\Delta_1=\log\coth(\beta J)\)。配分函数 \(Z_{1,M}=\Lambda_0^M+\Lambda_1^M\)，自由能 \(-\beta f_1=\log(2\cosh\beta J)+o(1)\) 与 Onsager 解的高温极限一致；纵向关联长度 \(\xi_\parallel=1/\log\coth(\beta J)\) 即一维 Ising 的精确值。

Perron--Frobenius 定理可推广至紧正算子（Krein--Rutman 定理）：若 \(\mathscr T\) 是 Banach 格上的紧正算子且谱半径 \(r(\mathscr T)>0\)，则 \(r(\mathscr T)\) 是具有正特征向量的单本征值。这把有限维 transfer matrix 的主本征值论证直接延拓到 \(N\to\infty\) 的柱面极限，是 Ruelle 证明 \(\log\Lambda_0(N)\) 凸性从而存在 thermodynamic limit 的关键工具。次主本征值在临界点处与主本征值等模，对应关联长度发散与相变。
\end{derivation}

\section{谱表示与 observable-dependent 相关长度}

transfer spectrum 不只决定自由能，也决定沿纵向的相关衰减。令 \(\widehat O\) 是由一行局域观测量定义的对角算符，取 \(M\to\infty\) 后，间隔 \(r\) 行的 connected correlation 可写成
\begin{equation}
 \langle O_0O_r\rangle_c
 =\sum_{a\ge1}
 C_a(O)
 \left(\frac{\Lambda_a}{\Lambda_0}\right)^r,
 \label{eq:phys-sm-transfer-correlation-spectrum}
\end{equation}
其中
\[
 C_a(O)=
 \langle0|\widehat O|a\rangle
 \langle a|\widehat O|0\rangle
\]
在 \(\mathsf T_N\) 对称实矩阵时非负。于是相关长度不是“第二大本征值”这一句话就完全决定，而要先看观测算符是否与该谱态有非零矩阵元。定义与 \(O\) 耦合的最低非零谱隙
\begin{equation}
 \xi_O^{-1}(N)
 =\min_{a:C_a(O)\neq0}
 \log\frac{\Lambda_0}{|\Lambda_a|}.
 \label{eq:phys-sm-observable-correlation-length}
\end{equation}
这条 selection rule 在后面区分单费米 transfer mass 与具体自旋相关长度时至关重要。

\begin{derivation}[transfer 谱到 connected correlation]
设 \(|0\rangle\) 是归一化 Perron--Frobenius 主本征矢，\(\mathsf T_N|0\rangle=\Lambda_0|0\rangle\)，并将 \(\mathsf T_N\) 的正交本征基记为 \(\{|a\rangle\}_{a=0}^{2^N-1}\)，满足
\[
 \mathsf T_N|a\rangle=\Lambda_a|a\rangle,
 \qquad
 \langle a|b\rangle=\delta_{ab}.
\]
对一行局域观测量 \(O\)，定义对角算符 \(\widehat O\)，并设两观测点纵向间隔 \(r\) 行。由 transfer 算子的物理意义，未归一化两点函数为
\[
 \langle0|\widehat O\mathsf T_N^r\widehat O|0\rangle
 \Lambda_0^{M-r-2},
\]
而配分函数主项为 \(\Lambda_0^M\)。在长柱面极限 \(M\to\infty\) 下取比值，所有 \(\Lambda_a/\Lambda_0\) 的 \(M\) 次幂中只有 \(a=0\) 项保留，归一化给出
\[
 \langle O_0O_r\rangle
 =\frac{\langle0|\widehat O\mathsf T_N^r
 \widehat O|0\rangle}{\Lambda_0^r}.
\]
将 \(\mathsf T_N^r\) 作用在本征基上，
\[
 \mathsf T_N^r|a\rangle=\Lambda_a^r|a\rangle,
\]
于是
\[
 \langle0|\widehat O\mathsf T_N^r\widehat O|0\rangle
 =\sum_a\Lambda_a^r
 \langle0|\widehat O|a\rangle
 \langle a|\widehat O|0\rangle,
\]
其中插入单位分解 \(\mathbb I=\sum_a|a\rangle\langle a|\) 于 \(\mathsf T_N^r\) 两侧。除以 \(\Lambda_0^r\) 后
\[
 \begin{aligned}
 \langle O_0O_r\rangle
 &=\sum_a
 \langle0|\widehat O|a\rangle
 \langle a|\widehat O|0\rangle
 \left(\frac{\Lambda_a}{\Lambda_0}\right)^r\\
 &=\langle0|\widehat O|0\rangle^2
 +\sum_{a\ge1}
 \langle0|\widehat O|a\rangle
 \langle a|\widehat O|0\rangle
 \left(\frac{\Lambda_a}{\Lambda_0}\right)^r.
 \end{aligned}
\]
由 \(|0\rangle\) 是 Perron--Frobenius 主本征矢，\(\langle0|\widehat O|0\rangle=\langle O\rangle\) 正是观测量 \(O\) 的热力学平均，因此 \(a=0\) 项为 \(\langle O\rangle^2\)。减去它便得 connected correlation
\[
 \langle O_0O_r\rangle_c
 :=\langle O_0O_r\rangle-\langle O\rangle^2
 =\sum_{a\ge1}
 C_a(O)
 \left(\frac{\Lambda_a}{\Lambda_0}\right)^r,
\]
其中 \(C_a(O)=\langle0|\widehat O|a\rangle\langle a|\widehat O|0\rangle\)，即\eqnrefs{eq:phys-sm-transfer-correlation-spectrum}。由 Perron--Frobenius 严格主本征值 \(|\Lambda_a|<\Lambda_0\)（\(a\ge1\)），每一项指数衰减。设 \(a^*\) 为使 \(C_a(O)\ne0\) 的最小 \(\log(\Lambda_0/|\Lambda_a|)\) 指标，则 \(r\to\infty\) 时主导项为 \(a^*\)，相关长度
\[
 \xi_O^{-1}
 =\log\frac{\Lambda_0}{|\Lambda_{a^*}|}
 =\min_{a:C_a(O)\ne0}
 \log\frac{\Lambda_0}{|\Lambda_a|},
\]
即\eqnrefs{eq:phys-sm-observable-correlation-length}。若 \(O\) 与所有非主谱态正交，则 connected correlation 恒为零，相关长度为无穷；selection rule 因此不是技术细节，而是观测量与谱的物理耦合条件。
\end{derivation}

\section{Pauli 表示与对偶耦合}
\label{sec:phys-sm-pauli}

在单点基 \(|+\rangle,|-\rangle\) 上定义 Pauli 算符 \(Z|\sigma\rangle=\sigma|\sigma\rangle\) 和翻转算符 \(X|\sigma\rangle=|-\sigma\rangle\)。行内水平因子写成
\begin{equation}
 \mathsf V_h
 =\exp\!\left(K_h\sum_{j=1}^{N}Z_jZ_{j+1}\right).
 \label{eq:phys-sm-horizontal-transfer-factor}
\end{equation}
单点垂直权重矩阵
\[
 \mathsf W_v=
 \begin{pmatrix}
 \ee^{K_v}&\ee^{-K_v}\\
 \ee^{-K_v}&\ee^{K_v}
 \end{pmatrix}
\]
可以写成 \(X\) 的指数。定义对偶耦合 \(K_v^*>0\) 使
\begin{equation}
 \tanh K_v^*=\ee^{-2K_v},
 \qquad
 \sinh(2K_v)\sinh(2K_v^*)=1,
 \label{eq:phys-sm-vertical-dual-coupling}
\end{equation}
则
\begin{equation}
 \mathsf W_v
 =\sqrt{2\sinh(2K_v)}\,\ee^{K_v^*X}.
 \label{eq:phys-sm-vertical-pauli-factor}
\end{equation}
因此
\begin{equation}
 \mathsf T_N
 =\left[2\sinh(2K_v)\right]^{N/2}
 \mathsf V_h^{1/2}\mathsf V_v\mathsf V_h^{1/2},
 \qquad
 \mathsf V_v=\exp\!\left(K_v^*\sum_{j=1}^{N}X_j\right).
 \label{eq:phys-sm-transfer-pauli-factorization}
\end{equation}

\begin{derivation}[垂直权重的指数表示]
单点垂直 Boltzmann 矩阵为
\[
 \mathsf W_v=
 \begin{pmatrix}
 \ee^{K_v}&\ee^{-K_v}\\
 \ee^{-K_v}&\ee^{K_v}
 \end{pmatrix},
\]
对角元与非对角元分别为 \(\ee^{K_v}\) 与 \(\ee^{-K_v}\)，二者之比为 \(\ee^{-2K_v}\)。由 \(X^2=\mathbb I\)，指数 \(\ee^{K_v^*X}\) 严格展开为
\[
 \begin{aligned}
 \ee^{K_v^*X}
 &=\sum_{n=0}^\infty\frac{(K_v^*)^nX^n}{n!}\\
 &=\sum_{m=0}^\infty\frac{(K_v^*)^{2m}}{(2m)!}\mathbb I
 +\sum_{m=0}^\infty\frac{(K_v^*)^{2m+1}}{(2m+1)!}X\\
 &=\cosh K_v^*\,\mathbb I
 +\sinh K_v^*\,X.
 \end{aligned}
\]
在基 \(\{|+\rangle,|-\rangle\}\) 下 \(X=\bigl(\begin{smallmatrix}0&1\\1&0\end{smallmatrix}\bigr)\)，因此
\[
 \ee^{K_v^*X}
 =\begin{pmatrix}
 \cosh K_v^*&\sinh K_v^*\\
 \sinh K_v^*&\cosh K_v^*
 \end{pmatrix},
\]
其对角元与非对角元之比为 \(\tanh K_v^*\)。令整体标量因子为 \(A_v>0\)，要求
\[
 A_v\,\ee^{K_v^*X}=\mathsf W_v,
\]
比较矩阵元得方程组
\[
 A_v\cosh K_v^*=\ee^{K_v},
 \qquad
 A_v\sinh K_v^*=\ee^{-K_v}.
\]
两式相除给出
\[
 \tanh K_v^*
 =\frac{\ee^{-K_v}}{\ee^{K_v}}
 =\ee^{-2K_v},
\]
即\eqnrefs{eq:phys-sm-vertical-dual-coupling}中第一式。为得第二式，将两式分别平方后相减，
\[
 \begin{aligned}
 A_v^2(\cosh^2 K_v^*-\sinh^2 K_v^*)
 &=\ee^{2K_v}-\ee^{-2K_v},
 \end{aligned}
\]
由双曲恒等式 \(\cosh^2 t-\sinh^2 t=1\) 与 \(\sinh(2K_v)=\tfrac12(\ee^{2K_v}-\ee^{-2K_v})\)，
\[
 A_v^2=2\sinh(2K_v).
\]
再由 \(\tanh K_v^*=\ee^{-2K_v}\) 解出
\[
 \sinh(2K_v^*)=\frac{2\tanh K_v^*}{1-\tanh^2 K_v^*}
 =\frac{2\ee^{-2K_v}}{1-\ee^{-4K_v}}
 =\frac{1}{\sinh(2K_v)},
\]
即 \(\sinh(2K_v)\sinh(2K_v^*)=1\)，亦即\eqnrefs{eq:phys-sm-vertical-dual-coupling}中第二式。综合 \(A_v=\sqrt{2\sinh(2K_v)}\) 与 \(K_v^*\) 的定义，
\[
 \mathsf W_v
 =\sqrt{2\sinh(2K_v)}\,\ee^{K_v^*X},
\]
即\eqnrefs{eq:phys-sm-vertical-pauli-factor}。把每个站点的垂直因子张量相乘得到 \(\mathsf V_v=\exp\!\bigl(K_v^*\sum_{j=1}^NX_j\bigr)\)，并在两侧各放一半水平权重 \(\mathsf V_h^{1/2}\) 以保持 \(\mathsf T_N\) 的对称分解，即得\eqnrefs{eq:phys-sm-transfer-pauli-factorization}。
\end{derivation}

\eqnrefs{eq:phys-sm-transfer-pauli-factorization} 仍是 \(2^N\) 维谱问题。平移对称只能按总晶格动量分块，并不能把指数维度降到线性规模。真正的可解结构是：一次自旋轴旋转后，\(\sum X_j\) 与 \(\sum Z_jZ_{j+1}\) 都成为 Jordan--Wigner 费米子的二次型，而二次 Clifford 算符的共轭作用只在线性 Majorana 空间中演化。因此必须先处理有限环边界项：它决定费米子奇偶与 NS/R 动量网格，不能在取热力学极限之前省略。

\begin{derivation}[Jordan--Wigner 变换与 Clifford 代数线性化]
Jordan--Winger 变换把 \(N\) 个 Pauli 自旋映射为 \(N\) 个费米子模式，使 transfer matrix 从 \(2^N\) 维一般矩阵降为 \(N\) 维自由费米子问题。分四步建立。

弦算符与费米子定义。对周期链 \(j=1,\ldots,N\)，定义弦算符（Jordan--Wigner 尾）
\begin{equation*}
  S_j=\prod_{k=1}^{j-1}Z_k,\qquad S_1=\mathbb I.
\end{equation*}
由 \(Z_k^2=\mathbb I\) 且不同位置 Pauli 对易，\(S_j^2=\mathbb I\)。定义 Majorana 算符
\begin{equation*}
  \gamma_j=S_j X_j,\qquad \tilde\gamma_j=S_j Y_j,\qquad j=1,\ldots,N.
\end{equation*}
由 \(X_jY_j=-Y_jX_j=\ii Z_j\) 和弦算符的对易性，直接验证 Clifford 关系
\begin{equation*}
  \{\gamma_j,\gamma_k\}=2\delta_{jk},\quad\{\tilde\gamma_j,\tilde\gamma_k\}=2\delta_{jk},\quad\{\gamma_j,\tilde\gamma_k\}=0.
\end{equation*}
复费米子 \(c_j=\frac{1}{2}(\gamma_j+\ii\tilde\gamma_j)\) 满足 \(\{c_j,c_k^\dagger\}=\delta_{jk}\)。

Pauli 二次型的费米子化。翻转算符 \(X_j=S_j\gamma_j\)（因 \(S_j^2=\mathbb I\)），但弦因子 \(S_j\) 使 \(\sum_j X_j\) 不是简单费米子二次型。关键技巧：做自旋轴旋转 \(X\to Y, Y\to-X\)（等价于 \(\gamma_j\to\tilde\gamma_j, \tilde\gamma_j\to-\gamma_j\)），使 transfer matrix 中两因子的弦因子结构互补。旋转后
\begin{equation*}
  \mathsf V_v\propto\exp\!\left(K_v^*\sum_j\tilde\gamma_j S_j\right),\qquad\mathsf V_h=\exp\!\left(K_h\sum_j\ii\gamma_j\tilde\gamma_{j+1}\cdot R_N\right),
\end{equation*}
其中 \(R_N=\prod_{k=1}^N Z_k=(-1)^{N_f}\) 是宇称算符（总费米子数奇偶）。

宇称分解与边界项。\(R_N\) 与所有 \(c_j,c_j^\dagger\) 对易（宇称守恒），因此 Fock 空间分解为偶宇称 \(\mathcal F_+\)（\(R_N=+1\)）与奇宇称 \(\mathcal F_-\)（\(R_N=-1\)）两个不变子空间。在每个子空间中 \(R_N\) 替换为常数 \(\pm1\)，\(\mathsf V_h\) 成为无弦因子的自由费米子二次型：
\begin{equation*}
  \mathsf V_h^{(\pm)}=\exp\!\left(\pm\ii K_h\sum_{j=1}^{N-1}\gamma_j\tilde\gamma_{j+1}\mp\ii K_h\gamma_N\tilde\gamma_1\right).
\end{equation*}
偶宇称给出周期（Ramond, R）边界 \(\gamma_{N+1}=\gamma_1\)，奇宇称给出反周期（Neveu--Schwarz, NS）边界 \(\gamma_{N+1}=-\gamma_1\)。这把 \(2^N\) 维问题分解为两个 \(2^{N-1}\) 维自由费米子问题。

Fourier 模式与 Bogoliubov 对角化。在每个宇称扇区，周期/反周期边界给出动量网格
\begin{equation*}
  k_m=\frac{2\pi m}{N}\quad(\text{R},\;m=0,1,\ldots,N-1),\qquad k_m=\frac{2\pi(m+1/2)}{N}\quad(\text{NS},\;m=0,1,\ldots,N-1).
\end{equation*}
Fourier 变换 \(c_j=\frac{1}{\sqrt N}\sum_k c_k\ee^{\ii kj}\) 后，每个动量模式 \((k,-k)\) 耦合为 \(2\times2\) 块，由 Bogoliubov 变换 \(\begin{pmatrix}\cos\theta_k&\sin\theta_k\\-\sin\theta_k&\cos\theta_k\end{pmatrix}\) 对角化，混合角 \(\tan(2\theta_k)=\frac{\sinh(2K_h)\sin k}{\cosh(2K_h)-\sinh(2K_h)\cos k}\)。本征值
\begin{equation}
  \epsilon(k)=\operatorname{arcosh}\!\left[\cosh(2K_h)\cosh(2K_v^*)-\sinh(2K_h)\sinh(2K_v^*)\cos k\right],
  \label{eq:phys-sm-fermion-dispersion}
\end{equation}
即 Onsager 色散关系。transfer matrix 本征值为各模式能量的乘积，自由能由 \(\frac{1}{N}\sum_k\epsilon(k)\to\int\frac{\dd k}{2\pi}\epsilon(k)\) 给出。临界条件 \(\epsilon(0)=0\)（R 扇区 \(k=0\) 模式软化）恢复\eqnrefs{eq:phys-sm-square-critical-k}。

Jordan--Wigner 弦 \(S_j=\prod_{k<j}Z_k\) 是一维有序化：它利用链的线性序把局域 Pauli 映射为局域费米子。二维中弦因子变为面填充算符，一般无法消除——这是二维 Ising 精确可解但三维不可解的根本原因。环面边界处弦因子绕回，产生 \(R_N\) 宇称分支，对应\eqnrefs{eq:phys-sm-torus-h1} 的两个同调生成元。这一联系把 Jordan--Wigner 边界项与曲面同调统一：弦绕行数 = 同调绕行数 = spin structure 选择。
\end{derivation}

上述二次费米子形式可用 Bogoliubov 变换直接求出谱隙。
自由费米子二次型由 Bogoliubov 正则变换对角化，给出精确色散和谱隙。分三步完成。

一般自由费米子 Hamiltonian \(H=\sum_{jk}(A_{jk}c_j^\dagger c_k+\frac12 B_{jk}c_j^\dagger c_k^\dagger+\text{h.c.})\)，\(A\) Hermite，\(B\) 反对称。Bogoliubov 变换 \(d_k=\sum_j(u_{kj}c_j+v_{kj}c_j^\dagger)\) 要求 \(d_k\) 满费米子关系且 \(H=\sum_k\epsilon_k d_k^\dagger d_k+\text{const}\)。代入得 Bogoliubov--de Gennes (BdG) 方程
\begin{equation*}
  \begin{pmatrix}A&B\\-B^*&-A^*\end{pmatrix}\begin{pmatrix}u_k\\v_k\end{pmatrix}=\epsilon_k\begin{pmatrix}u_k\\v_k\end{pmatrix},
\end{equation*}
\(2N\times2N\) 特征值问题，正负配对 \(\pm\epsilon_k\)。

对 Ising transfer matrix，Fourier 后 \(A_k=K_h\cos k+K_v^*\)，\(B_k=K_h\sin k\)（各向异性推广直截了当）。BdG 矩阵
\begin{equation*}
  \begin{pmatrix}K_h\cos k+K_v^*&K_h\sin k\\-K_h\sin k&-(K_h\cos k+K_v^*)\end{pmatrix},
\end{equation*}
本征值 \(\epsilon(k)=\pm\sqrt{(K_h\cos k+K_v^*)^2+(K_h\sin k)^2}=\pm\sqrt{K_h^2+K_v^{*2}+2K_hK_v^*\cos k}\)。这与\eqnrefs{eq:phys-sm-fermion-dispersion} 在参数化 \(\cosh\epsilon=\cosh(2K_h)\cosh(2K_v^*)-\sinh(2K_h)\sinh(2K_v^*)\cos k\) 下一致（差一个整体重标，来自 transfer matrix 与 Hamiltonian 的指数关系）。

谱隙 \(\Delta=\min_k\epsilon(k)=\epsilon(0)\)（R 扇区）。由\eqnrefs{eq:phys-sm-fermion-dispersion}，
\begin{equation*}
  \cosh\Delta=\cosh(2K_h)\cosh(2K_v^*)-\sinh(2K_h)\sinh(2K_v^*)=\cosh(2K_h-2K_v^*).
\end{equation*}
各向同性 \(K_h=K_v=K\)，\(K_v^*=K^*\)（对偶），\(\Delta=2|K-K^*|\)。临界点 \(K=K_c\)（即 \(K=K^*\)）时 \(\Delta=0\)，谱隙闭合，关联长度 \(\xi=1/\Delta\to\infty\)。临界附近 \(K=K_c+\delta\)，\(K^*=K_c-\delta\)（对偶性），\(\Delta\approx4\delta\)，\(\xi\approx1/(4\delta)\)，给出 \(\nu=1\)。这一结果独立于 Kramers--Wannier 对偶，纯粹来自 Bogoliubov 谱分析，是对偶论证的严格确认。

有限环 \(N\) 上 NS 扇区最低动量 \(k_{\min}=\pi/N\)，有限尺寸谱隙 \(\Delta_N\approx\Delta+\frac{\pi^2}{2N^2\epsilon(0)}\)（小 \(k\) 展开）。临界点 \(\Delta=0\) 时 \(\Delta_N\sim\pi/N\)，对应 central charge \(c=1/2\) 的 Ising CFT：有限尺寸自由能修正 \(f_N=f_\infty-\frac{\pi c}{6N^2}\)，\(c=1/2\) 由三个 primary fields \((1,\sigma,\epsilon)\) 的 conformal weights \((0,1/16,1/2)\) 确定。这把 transfer matrix 谱与共形场论直接联系起来。

谱隙 \(\Delta\) 是关联长度倒数 \(\xi^{-1}\) 的量级，控制所有关联函数的指数衰减率。临界点 \(\Delta=0\) 意味着关联长度发散，系统出现长程序（对 \(T<T_c\)）或代数衰减关联（对 \(T=T_c\)）。Bogoliubov 方法给出的 \(\nu=1\) 与 Onsager 精确解一致，是自由费米子可积性的直接体现——相互作用系统通常无此精确谱隙。

Bogoliubov 变换 \((c,c^\dagger)\mapsto(d,d^\dagger)\) 保持费米子反对易关系，是辛群 \(\mathrm{Sp}(2N,\mathbb R)\cap\mathrm O(2N)=\mathrm U(N)\) 的元素。BdG 矩阵 \(\begin{pmatrix}A&B\\-B^*&-A^*\end{pmatrix}\) 属于 \(\mathfrak{so}(2N)\)（反 Hermite 化后），其本征值实部给出色散。这一辛结构保证了正负本征值配对 \(\pm\epsilon_k\)（粒子-空穴对称），是自由费米子可积性的群论基础。相互作用破坏辛结构（BdG 不再封闭），使精确对角化不可能。

最近邻相互作用 Ising（\(J_1\sigma_i\sigma_{i+1}+J_2\sigma_i\sigma_{i+2}\)）在 \(J_2/J_1\) 一般时不可积：Jordan--Wigner 后 Hamiltonian 含四费米子项 \(\sim c^\dagger c^\dagger cc\)，Bogoliubov 变换无法消去。例外是 \(J_2=0\)（自由费米子）和 \(J_2/J_1=\pm1/2\)（Baxter 可积点），后者有额外的 Yang--Baxter 结构保证可积性。一般相互作用系统的谱只能数值或重整化群处理，精确谱隙不存在——这把 Ising 精确解的"特殊性"精确界定：自由费米子（辛对称性）或 Yang--Baxter（量子群对称性）是可积的两大机制。

SK 模型（Sherrington--Kirkpatrick）的无序平均使谱问题变为随机矩阵本征值问题（Wigner 半圆律），与 Ising 的确定性 Bogoliubov 谱截然不同。SK 的临界行为由 Parisi 超度量结构描述，关联长度发散机制是 replica 对称破缺而非谱隙闭合。这对比凸显了 Ising 精确可解性的独特性：确定性格点 + 自由费米子 \(\to\) Bogoliubov 精确谱；无序 + 平均场 \(\to\) 随机矩阵 + 超度量。

转移矩阵还可通过 Suzuki--Trotter 分解与一维量子链联系起来。
Suzuki--Trotter 分解把 \(d\) 维经典统计模型的 transfer matrix 映射为 \((d-1)\) 维量子模型的虚时演化，建立量子--经典对应。分四步建立。

对算符 \(A,B\)，\(\ee^{A+B}=\lim_{n\to\infty}(\ee^{A/n}\ee^{B/n})^n\)。误差估计：\(\ee^{A/n}\ee^{B/n}=\ee^{(A+B)/n+[A,B]/(2n^2)+O(n^{-3})}\)，故 \(n\) 步分解误差 \(O(1/n)\)。对 Ising transfer matrix \(\mathsf T=\mathsf V_h\mathsf V_v\)，Trotter 分解把非对易乘积化为交替指数的乘积，每步可显式计算。

2D 经典 Ising 的 transfer matrix \(\mathsf T_N=\mathsf V_h^{1/2}\mathsf V_v\mathsf V_h^{1/2}\)（\eqnrefs{eq:phys-sm-transfer-pauli-factorization}）。连续极限 \(M\to\infty\)（纵向层数），\(\mathsf T_N^M=\ee^{-M H_{\rm Q}}\)，量子 Hamiltonian
\begin{equation*}
  H_{\rm Q}=-\sum_j(K_v^*\tau_j^x+K_h\tau_j^z\tau_{j+1}^z)+\text{const},
\end{equation*}
即 1D 横场量子 Ising 链。经典 2D Ising 的临界点 \(K_c\) 对应量子 1D Ising 的量子临界点 \(K_v^*/K_h=g_c\)（横场/纵向耦合比）。经典关联长度 \(\xi_{\rm class}\) 对应量子关联长度 \(\xi_{\rm quant}=1/\Delta\)（\(\Delta\) 是量子谱隙）。

量子--经典对应使临界指数按 \(d_{\rm class}=d_{\rm quant}+1\) 平移：经典 \(d\) 维的 \(\nu_{\rm class}\) 对应量子 \((d-1)\) 维的 \(\nu_{\rm quant}\)。2D 经典 Ising（\(\nu=1\)）对应 1D 量子 Ising（\(\nu=1\)，不位移因 \(\nu\) 在量子中定义不同）。更精确地，量子相变由 gap \(\Delta\sim|g-g_c|^{\nu z}\)（\(z\) 是 dynamical exponent）控制，经典对应 \(z=1\)（Lorentz 不变临界点）。2D 经典 Ising 的 \(\nu=1\)，\(z=1\)，对应 1D 量子 Ising 的 \(\nu z=1\)。

\begin{itemize}
 \item[(i)] 量子相变：1D 横场 Ising 链在 \(g=g_c\) 有量子相变，从有序相（\(g<g_c\)，\(\langle\tau^z\rangle\ne0\)）到无序相（\(g>g_c\)，\(\langle\tau^z\rangle=0\)）。
 \item[(ii)] 绝热量子计算：Kitaev 的 honeycomb 模型通过 Suzuki--Trotter 映射到 2D Ising，可积性来自同一 Jordan--Wigner/Clifford 结构。
 \item[(iii)] 蒙特卡洛：Suzuki--Trotter 分解把量子统计力学化为经典配分函数（世界线表示），使量子 Monte Carlo 成为可能——连续极限误差 \(O(\Delta\tau^2)\) 由分解阶数控制。
 \item[(iv)] AdS/CFT 对应：量子--经典对应的推广，\(d\) 维量子场论映射到 \((d+1)\) 维经典引力，Ising/CFT 对应是其离散版本。
\end{itemize}
这一对应把 Ising 精确解从"经典统计力学技巧"提升为"量子--经典对偶的原型"。
